\chapter{Introduction}
\label{chap:introduction}

\section{Background}
\label{sec:background}

Music, podcasts and audiobooks can form part of the same person's listening
activities, while being accessed through different applications. In such a
situation, finding audio involves choosing both what to listen to and where
to look for it. A listener may save an item in one application, continue a
different item elsewhere and later need to remember which application holds
each one. These possible organisational difficulties motivate the project.

Finding content is only one part of this problem. Saving an interesting item
should also make it possible to find that item again. Similarly, returning to
an unfinished podcast or audiobook requires a recognisable route back to the
content and its listening position. For example, a student who stops listening
to an audiobook during a study break should be able to identify the unfinished
book and continue it later.

Combining the three audio types could provide a common place for these
activities, but it also raises a design question. A mixed collection may be
difficult to navigate if its content types are unclear or if too many items
compete for attention. Music tracks, podcast episodes and audiobook chapters
also need to remain identifiable within a shared interface. Common navigation
and saving controls therefore have to be balanced with clear categories and
information about each item.

Omnisound is the project's response to this question: a web-interface
prototype for finding, saving and resuming audio in one organised place. Its
concept sheet identifies students who regularly use music, podcast or
audiobook applications as the intended users, including listeners who use
several content types and those who mainly listen to music
\cite[p.~1]{omnisoundconcept}. The concept treats losing track of saved or
unfinished content as a working problem to investigate through the design
and evaluation of the interface.

The high-fidelity prototype is designed primarily for a desktop browser. It
provides a shared search, content-type filters, a library of saved items and a
Continue Listening area. Its two documented task flows are finding and saving
an item, then locating it in the library, and resuming an unfinished podcast
or audiobook \cite[Sec.~``Testable user flows'']{omnisoundv1}. These flows give
the project a specific interaction focus within the broader audio concept.

\section{Purpose}
\label{sec:purpose}

The purpose of the project is to investigate how a unified audio interface
can support finding, organising and resuming different types of audio content
while keeping the experience clear and manageable for the intended users.
The motivation is to explore whether a common search and library can reduce
the effort of returning to wanted content without making the differences
between music, podcasts and audiobooks harder to understand.

This investigation also provides an opportunity to apply user-centred design
within the HCI course. The course project requires an iterative process and
usability testing in at least one iteration \cite[pp.~2--3]{p3project}.
For Omnisound, the purpose is therefore pursued through requirements,
alternative designs, prototypes and evaluation. The desired qualities are
ease of learning, clear organisation and useful feedback; whether the design
supports these qualities must be assessed through the project work.

\section{Goals}
\label{sec:goals}

The completed project should provide an identified set of user needs and a
prioritised set of interface requirements. The requirements should make clear
which listening activities the interface is intended to support and how they
relate to the information collected from the target group. This provides a
basis for assessing the design beyond its visual appearance.

The design deliverables should include documented alternatives, low-fidelity
sketches and interactive high-fidelity prototypes. Alternatives should cover
the organisation of audio categories and the presentation of content, such as
category tabs compared with a mixed view with filters, and cards compared
with lists. The sketches should preserve the early ideas, while the V1 and V2
prototypes should make the selected task flows available for examination.
Together, these materials should show how the interface developed and why
particular solutions were chosen.

A further deliverable is a documented usability evaluation of the important
tasks: finding and saving content, locating saved content and returning to
unfinished spoken audio. The evaluation should provide observations of task
completion and difficulties, together with an account of how those
observations informed design decisions. A record of design iterations and
reasoned conclusions should connect the requirements, prototypes and
evaluation. These goals concern identifiable project outputs; they do not
assume that all proposed changes will prove effective.

\section{Delimitations}
\label{sec:delimitations}

Omnisound is an interface prototype. The project focuses on navigation,
content organisation and the visibility of actions and listening progress.
It does not implement music-streaming infrastructure, integration with
commercial audio services, payment or subscription processing, licensing
systems, backend databases or production account management. This scope is
consistent with the prototype documentation, which describes a front-end-only
implementation \cite[opening description]{omnisoundv1}. Limiting the
implementation in this way allows the design work to concentrate on the
selected interactions within the course project.

Search uses a sample catalogue, and the player controls and displayed progress
represent the intended listening experience. They are sufficient to explore
how users move through the interface, recognise content and interpret
feedback. They do not demonstrate reliable audio delivery, synchronisation
between devices or the operation of a real streaming service. Conclusions
must therefore concern interaction with the prototype rather than the
technical performance of a complete product.

The task scope is also limited. Search and save, retrieval from the library
and continuation of unfinished podcasts or audiobooks receive priority.
Other possible activities, such as managing a production account or exploring
a large commercial catalogue, are outside this evaluation focus. The primary
context is desktop use, so the project cannot establish usability across all
mobile devices or listening environments. Likewise, conclusions drawn from
evaluation with the intended student group should be considered in relation
to that group and the tasks tested.

Following the first usability test, a second version of Omnisound was
developed to address the problems identified during evaluation
\cite[pp.~15--18]{omnisoundtestv1}. The changes focused on reducing the size
of the Home banner, making Continue Listening easier to discover, and
improving the clarity of the Save/Saved control
\cite[Sec.~``Iteration 2 redesign'']{omnisoundv2}. A second usability test was
then conducted using V2 and the same three core tasks, allowing the group to
compare the two iterations \cite[pp.~3--4, 13]{omnisoundtestv2}. The testing
procedures and findings are presented in the Method and Results chapters.
Since Omnisound remains an interface prototype, the conclusions concern the
tested interactions rather than the technical performance of a complete
streaming service.

\section{Division of Work}
\label{sec:division-of-work}

The project group consists of Obid Alrhman Aljabi, Yacoub Dawli and Sarshar
Chavoshi Rad. The Project Session 2 data collection plan assigns Yacoub Dawli
the interviewer role and assigns Obid Alrhman Aljabi and Sarshar Chavoshi Rad
the recording of handwritten interview notes \cite[p.~2]{interviewplan}.
This documents the planned allocation of those roles. It does not confirm
which interviews were completed or who carried out the later design,
implementation and evaluation work.

\noindent\textbf{TODO -- group confirmation:} Record each member's actual
contributions to data collection and analysis, sketches, V1 and V2,
usability-test planning and execution, and report writing. Confirm whether the
planned interview roles were followed, and identify the report sections
written or revised by each member before submission.
